\begin{claim}{D.52}{Proved}{An abstract allowed list for every finite full torus subgroup}
Let $T$ be the full once-punctured-torus loop-outer
group of D.50. Every finite subgroup of $T$ has
one of the following twelve abstract types:
\[
\begin{gathered}
1,\ C_2,\ C_3,\ C_4,\ C_2\times C_2,\ C_6,\ S_3,\ D_8,\\
C_{12},\ D_{12},\ \operatorname{Dic}_{12},\ D_{24}.
\end{gathered}                                             \tag{D.52a}
\]
Here $D_{2n}$ has order $2n$, and
$\operatorname{Dic}_{12}=C_3\rtimes C_4$, with
a generator of $C_4$ inverting $C_3$.
The nine types obtained by omitting
$C_{12},\operatorname{Dic}_{12},D_{24}$ all occur.
Occurrence of those three remaining types is
not determined by this allowed list.
The subsequent full covering arguments in D.54--D.55
exclude all three.

Every subgroup of type $S_3,D_{12}$ or $D_{24}$
has the global usual reflection multiplier:
rotations have multiplier $1$ and reflections
have multiplier $-1$.
In a subgroup of type $\operatorname{Dic}_{12}$,
the subgroup of order three has multiplier $1$,
and every generator $b$ of a cyclic subgroup
of order four satisfies
\[
       \lambda(b)^2=1,\qquad
       \lambda(b)_2=1,\qquad\lambda(b)_3=-1.              \tag{D.52b}
\]
For a subgroup of type $C_{12}$, its order-four
generator instead has $\lambda(b)_2=\lambda(b)_3=1$,
and still has $\lambda(b)^2=1$.
These restrictions retain the unresolved cyclic
order-three and order-four multiplier questions.
\end{claim}
\begin{proof}
\emph{The two primary representations and their actual kernels.}
For a finite subgroup $G\leq T$, D.33 gives
\[
 |G|\mid24,\qquad
 A_p:G\longrightarrow\GL_2(\Z_p),\qquad
 \det A_p(g)=\lambda(g)_p.                              \tag{D.52c}
\]
It makes $A_p$ injective on every finite
$p$-subgroup. Put $K=\ker(A_2|_G)$.
An element of order two in $K$ would contradict
that injection. Cauchy's theorem and (D.52c)
therefore give
\[
                      K=1\quad\hbox{or}\quad K\cong C_3.
                                                               \tag{D.52d}
\]
We classify the actual image $H=A_2(G)$ and
then the extension by this specified kernel.
We do not identify $G$ with a representation
at one prime when the kernel is nontrivial.

\emph{Finite images over $\Q_2$.}
The two-subgroups of $H$ have the types in D.50.
Indeed a Sylow two-subgroup of the preimage
of any such subgroup maps isomorphically onto
it, since the kernel has odd order.
The matrix arguments in D.50 also explicitly
exclude the quaternion group of order eight.

Suppose $3$ divides $|H|$.
Its number of Sylow three-subgroups is either
one or four. We exclude four.
Conjugation gives a transitive homomorphism
$H\to S_4$. Its kernel has order at most six.
If that kernel had order divisible by three,
its unique Sylow three-subgroup would be
characteristic in the kernel and normal in $H$,
contrary to the assumed four Sylow subgroups.
Thus the transitive image has order divisible
by both four and three, and is $A_4$ or $S_4$.
The kernel consequently has order one or two.

There is no faithful two-dimensional
representation of $A_4$ over $\Q_2$.
Its normal Klein four subgroup would be
simultaneously diagonalizable and, by
faithfulness, would be the full group
of diagonal sign matrices. One of its
three involutions would therefore be $-I$.
Conjugation fixes this scalar, whereas the
order-three elements of $A_4$ cyclically
permute all three involutions.
This is a contradiction.
It excludes the cases with trivial kernel,
whose images are $A_4$ or $S_4$.
In the remaining case $|H|=24$ and the image
is $A_4$ with a kernel of order two.
The preimage $P$ of its Klein four subgroup
is a normal subgroup of order eight, hence
is $D_8$ by D.50. But
$|\operatorname{Aut}(D_8)|=8$: a generator
of its rotation subgroup has two possible
images, and a chosen reflection has four.
Every element of order three therefore
centralizes $P$. Multiplication by a Sylow
three-subgroup gives $H=P\times C_3$, which
has only one Sylow three-subgroup.
This contradiction proves that $H$ has a
normal subgroup $P_3$ of order three.

Let $R$ generate $P_3$.
Its minimal polynomial is
$X^2+X+1$, irreducible over $\Q_2$ by
reduction modulo two.
The basis $(v,Rv)$ for any nonzero vector $v$
identifies the matrix centralizer with
\[
 L_2^\times,\qquad L_2=\Q_2(\zeta_3).
                                                               \tag{D.52e}
\]
Under this identification, a commuting
endomorphism is multiplication by $a+b\zeta_3$.
The full normalizer consists exactly of these
multiplication maps and their compositions
with $\sigma(\zeta_3)=\zeta_3^{-1}$:
conjugation must send $R$ to $R$ or $R^{-1}$,
and division by $\sigma$ in the second case
gives a commuting map.

The field $L_2$ contains no element of order
four. To check this directly, the equation
$(a+b\zeta_3)^2=-1$ has coefficient equations
\[
                   a^2-b^2=-1,\qquad b(2a-b)=0.
                                                               \tag{D.52f}
\]
They require either $a^2=-1$ or $a^2=1/3$
in $\Q_2$. Both are impossible: these are
units congruent respectively to seven and
three modulo eight, whereas an odd square
is congruent to one.
Since $|H|\mid24$, it follows that
$C=H\cap L_2^\times$ is $C_3$ or $C_6$.
Here its two-part is contained in
$\{1,-1\}$, and its three-part is $P_3$,
so no general classification of local
roots of unity is required.

If $H=C$, these are the two possible groups.
Otherwise choose $B=a\sigma\in H\setminus C$.
Then
\[
 B^2=N_{L_2/\Q_2}(a)\,I
        \in C\cap\Q_2^\times=\{I,-I\}\cap C.             \tag{D.52g}
\]
The last scalar restriction follows because
a finite-order scalar of order dividing
twenty-four in $\Q_2$ is a sign: there are
no primitive cube or fourth roots in $\Q_2$.
For $C=C_3$, equation (D.52g) gives $S_3$.
For $C=C_6$, its two values give respectively
$D_{12}$ and $\operatorname{Dic}_{12}$.
Thus the images with order divisible by
three have the complete allowed list
\[
                  C_3,\ C_6,\ S_3,\ D_{12},
                             \operatorname{Dic}_{12}.  \tag{D.52h}
\]

\emph{The three-adic normalizer and its determinant.}
We will need a statement about any normal
subgroup $P_3\cong C_3$ of the actual $G$.
The map $A_3$ is injective on $P_3$.
Its generator has minimal polynomial
$X^2+X+1$ over $\Q_3$, irreducible since
its discriminant $-3$ has odd valuation.
The same centralizer and normalizer
calculation identifies
\[
 C_{\GL_2(\Q_3)}(A_3(P_3))=L_3^\times,\qquad
 N_{\GL_2(\Q_3)}(A_3(P_3))=L_3^\times\rtimes\langle\sigma\rangle,
 \quad L_3=\Q_3(\zeta_3).                              \tag{D.52i}
\]
This field also has no element of order four:
(D.52f) would require $a^2=-1$, impossible
modulo three, or $a^2=1/3$, impossible by
valuation. Thus every finite-order element
in its centralizer with order dividing
twenty-four lies in $\mu_6$ and has norm one.

There is also no element of norm $-1$ in $L_3$.
Here is an explicit calculation, including
possible nonintegral coordinates.
The norm of $a+b\zeta_3$ is
$a^2-ab+b^2$. Put
$m=\min(v_3(a),v_3(b))$ and write
$a=3^m u,b=3^m v$, where $u,v\in\Z_3$
and at least one is a unit.
When $u+v\not\equiv0\pmod3$, the norm of
$u+v\zeta_3$ has valuation zero and residue
$(u+v)^2=1$ modulo three.
When $u+v\equiv0\pmod3$, both are units;
writing $u=-v+3k$ gives
\[
                  u^2-uv+v^2
                       =3(v^2-3vk+3k^2),               \tag{D.52j}
\]
of valuation one.
The norm $-1$ would have valuation zero.
The second case cannot give that valuation,
since $2m+1\ne0$.
The first forces $m=0$ and residue one,
whereas $-1$ has residue two.
This proves the asserted norm obstruction.

An element of the finite image $A_3(G)$
outside the centralizer has the form
$B=a\sigma$, and $B^2=N(a)I$ is a
finite-order scalar, hence a sign.
The norm obstruction forces $B^2=I$.
In the basis $(1,\zeta_3)$, the determinant
of $\sigma$ is $-1$, so
$\det B=-N(a)=-1$.
Together with the centralizer calculation,
this proves the exact character identity
\[
 \lambda(g)_3=
 \begin{cases}
  1,&g\text{ centralizes }P_3,\\
 -1,&g\text{ inverts }P_3.
 \end{cases}                                            \tag{D.52k}
\]
The argument uses faithfulness only on
$P_3$, not on all of $G$.

\emph{Reconstructing the actual finite subgroup.}
If $K=1$ in (D.52d), the injection into
$H$ gives precisely D.50 and (D.52h).
Now let $K=P_3\cong C_3$.
A Sylow two-subgroup $Q$ of $G$ has order
$|G|/3$ and maps isomorphically onto $H$.
The actual multiplication map
\[
                  P_3\rtimes Q\longrightarrow G,
                         \qquad (a,q)\longmapsto aq,    \tag{D.52l}
\]
is an isomorphism, with action given by
conjugation. D.50 lists the five possibilities
for $Q$.
For $Q=1,C_2,C_4$, their homomorphisms
to $\operatorname{Aut}(C_3)=C_2$ give
$C_3$; $C_6$ or $S_3$; and $C_{12}$ or
$\operatorname{Dic}_{12}$, respectively.

For $Q=C_2\times C_2$ or $D_8$, D.50
already determines its global multiplier.
Equation (D.52k) therefore makes its
action on $P_3$ exactly the reflection
character of D.50.
Write $a$ for a generator of $P_3$.
For $Q=C_2\times C_2$, write $z$ for its
positive involution and $s$ for a negative
one. The element $az$ has order six,
is inverted by $s$, and generates $a,z$
through $(az)^4=a$ and $(az)^3=z$.
This gives an explicit $D_{12}$ presentation.
For $Q=D_8=\langle r,s\rangle$, where
$r$ is its order-four rotation, the
element $ar$ has order twelve, is inverted
by $s$, and satisfies
\[
                       (ar)^4=a,\qquad (ar)^9=r.
                                                               \tag{D.52m}
\]
It gives the explicit $D_{24}$ presentation.
These actual maps finish the allowed list.

\emph{Multipliers and realized cases.}
In every group just obtained that inverts
its normal $C_3$, the multiplier on that
$C_3$ is one: a conjugation-invariant
homomorphism into the abelian group
$\widehat\Z^\times$ takes an inverted
element to its own inverse, and a value
of order dividing three with square one
is one.
In $S_3,D_{12},D_{24}$ each reflection
inverts $C_3$ and hence has three-adic
multiplier $-1$ by (D.52k).
D.43 makes its full multiplier $-1$.
Every rotation is a product of two
reflections, proving the stated global
reflection character.
For $\operatorname{Dic}_{12}$, an
order-four generator inverts $C_3$,
so (D.52k) gives its negative
three-coordinate. D.50 gives its positive
two-coordinate and global square one.
For $C_{12}$ the same generator centralizes
$C_3$, giving instead the positive
three-coordinate. This proves all the
claimed multiplier restrictions.

D.50 realizes its five types by explicit
full profinite free-group automorphisms.
The other four realized types have the
following simultaneous exact surface model.
On $\R^2/\Z^2$ the integral matrices
\[
       R=\begin{pmatrix}0&-1\\1&1\end{pmatrix},
       \qquad
       S=\begin{pmatrix}0&1\\1&0\end{pmatrix}            \tag{D.52n}
\]
fix the origin and satisfy
$R^6=S^2=I$, $SRS=R^{-1}$.
Their twelve matrices $R^iS^j$ are distinct.
Removing the origin therefore gives an
actual $D_{12}$ action on a once-punctured
torus. Passing to its fundamental group,
then to its full profinite completion and
outer automorphisms, gives a homomorphism
to $T$. A small peripheral circle is
carried to an ellipse of local degree
$\det R=1$ or $\det S=-1$, so this
homomorphism preserves the stated full
cusp with the indicated multiplier.
The induced homology matrices are exactly
$R,S$ in $\GL_2(\widehat\Z)$; their
distinctness proves that the homomorphism
is injective.
The subgroups
$\langle R^2\rangle,\langle R\rangle,
\langle R^2,S\rangle$, and
$\langle R,S\rangle$ realize
$C_3,C_6,S_3,D_{12}$.

The program
\path{code/part_c/torus_finite_types_d52.py}
checks the matrix relations, the local
norm and square residues, the eight
automorphisms of $D_8$, and the explicit
finite semidirect-product identifications
in (D.52l)--(D.52m).
The prime-wise kernel and valuation
arguments above prove the full allowed
list. No assertion of occurrence or
nonoccurrence of its three remaining
types is inferred from those finite checks.
\end{proof}
